\section{Why Graph Cycles Are the Wrong Test}
\label{sec:bare-graph-not-enough}

For a finite abstract rewriting system the reduction graph \emph{is} the system: its vertices are the states and its directed edges are the one-step rewrites. Every property in \Cref{def:basic-notions}, confluence included, is therefore a property of that directed graph, and an algorithm can decide it from the graph by computing reachability. The question of this section is narrower. Graph theory and discrete geometry offer ready-made notions of flatness built from \emph{cycles}, and one may hope that one of them captures confluence. We show that the two most natural ones do not.

\begin{enumerate}
    \item \emph{Directed-cycle flatness} calls the reduction graph flat when it has no directed cycle.
    \item \emph{Undirected cycle rank} measures the dimension of the cycle space of the underlying simple undirected graph, i.e.\ the number of independent undirected cycles; rank zero means the graph is a forest.
\end{enumerate}

\begin{figure}[t]
\centering
\begin{tikzpicture}
  \begin{scope}
    \node[st] (A) at (0,0) {$A$};
    \node[st] (B) at (-0.9,-1.0) {$B$};
    \node[st] (C) at (0.9,-1.0) {$C$};
    \node[st] (D) at (0,-2.0) {$D$};
    \draw[rw] (A)--(B); \draw[rw] (A)--(C); \draw[rw] (B)--(D); \draw[rw] (C)--(D);
    \node[font=\small, align=center] at (0,-2.9) {(a) diamond $E001$};
    \node[font=\scriptsize, align=center] at (0,-3.6) {acyclic, cycle rank $1$\\ confluent};
  \end{scope}
  \begin{scope}[xshift=4.6cm]
    \node[st] (A) at (0,0) {$A$};
    \node[st] (B) at (-0.9,-1.0) {$B$};
    \node[nf] (C) at (0.9,-1.0) {$C$};
    \node[nf] (D) at (-0.9,-2.0) {$D$};
    \draw[rw] (A)--(B); \draw[rw] (A)--(C); \draw[rw] (B)--(D);
    \draw[defect] (D) .. controls (0.3,-2.2) and (0.9,-1.8) .. (C);
    \node[font=\small, align=center] at (0,-2.9) {(b) fork $E002$};
    \node[font=\scriptsize, align=center] at (0,-3.6) {acyclic, cycle rank $0$\\ not confluent};
  \end{scope}
  \begin{scope}[xshift=9.2cm]
    \node[st, minimum size=14pt] (s) at (0,-1.2) {$s$};
    \draw[rw] (s.north east) .. controls (0.9,0.2) and (-0.9,0.2) .. (s.north west);
    \node[font=\small, align=center] at (0,-2.9) {(c) one self-loop};
    \node[font=\scriptsize, align=center] at (0,-3.6) {directed cycle\\ confluent};
  \end{scope}
\end{tikzpicture}
\caption{Cycle structure says nothing about confluence. (a) The diamond is confluent and has an undirected cycle. (b) The fork has no cycles of either kind but is not confluent: its two normal forms $C$ and $D$ (boxed) have no common descendant, as indicated by the dotted line. (c) A single state with a self-loop is confluent, since every state reachable from $s$ is $s$ itself, yet it has a directed cycle.}
\label{fig:bare-graph-failure}
\end{figure}

\begin{table}[t]
    \caption{Cycle-based proxies on the curated finite systems $E001$--$E004$. All four graphs are acyclic, so directed-cycle flatness calls all of them flat, yet $E002$ and $E004$ are not confluent. Undirected cycle rank is positive on the confluent systems $E001$ and $E003$ and zero on the nonconfluent ones. The column of nonjoinable peaks is the local quantity that does track confluence (\Cref{sec:core-theorem}).}
    \label{tab:bare-graph-flatness-mismatch}
    \centering
    \smallskip
    \small
\begin{tabular}{@{}lcccc@{}}
\toprule
Example & \shortstack{Directed-cycle\\flat?} & \shortstack{Undirected\\cycle rank} & \shortstack{Nonjoinable\\peaks} & Confluent? \\
\midrule
$E001$ (diamond)           & yes & 1 & 0 & yes \\
$E002$ (fork)              & yes & 0 & 1 & no \\
$E003$ (delayed join)      & yes & 1 & 0 & yes \\
$E004$ (double fork)       & yes & 0 & 2 & no \\
\bottomrule
\end{tabular}

\end{table}

\begin{proposition}[Cycle-based proxies do not detect confluence]
\label{prop:bare-graph-not-diagnostic}
Among finite abstract rewriting systems, neither directed-cycle flatness nor vanishing undirected cycle rank implies confluence, and neither is implied by confluence. Directed-cycle flatness fails to imply confluence even among terminating systems, where it holds by definition.
\end{proposition}

\begin{proof}
All four claims are witnessed by the systems in \Cref{fig:bare-graph-failure}. The fork (b), with edges $A \to B$, $A \to C$, $B \to D$, has no directed cycle and its undirected graph is a tree, yet $C$ and $D$ are distinct normal forms reachable from $A$, so it is not confluent. The self-loop (c) has a directed cycle and is confluent. The diamond (a) has undirected cycle rank $1$ and is confluent, so confluence does not force rank zero. Finally, every terminating finite system is acyclic and hence directed-cycle flat, while the fork is terminating and not confluent.
\end{proof}

The exhaustive search of \Cref{subsec:boundary-search} finds these witnesses automatically: the smallest system on which directed-cycle flatness holds without confluence is the three-state fork $s_0 \to s_1$, $s_0 \to s_2$, and the smallest confluent system with a directed cycle is the one-state self-loop.

Why do cycle-based notions miss confluence? Because confluence compares pairs of rewrite sequences out of a common state, and locally, for instance in terminating systems where local confluence suffices, it is tested on \emph{peaks}, pairs of one-step rewrites $b \leftarrow a \to c$. A divergence need not contain a cycle, and the existence of a cycle does not determine whether its branches are joinable. In the fork the peak at $A$ cannot be joined, but nothing in the graph closes up around it; in the diamond the peak can be joined, and only then does it become part of a closed walk $A \to B \to D \leftarrow C \leftarrow A$. Cycles that do exist in a reduction graph may have nothing to do with peaks, as in the self-loop. A notion of flatness adapted to confluence must therefore first say \emph{which pairs of paths are being compared}. The reduction graph contains all the information needed to answer the question, but it does not single out the peaks as the objects on which the comparison takes place. The next section makes that comparison data explicit.
